\begin{abstract}
Biological knowledge bases such as the Gene Ontology (GO) annotation corpus underpin enrichment
analysis, function prediction and the training of machine-learning models, yet their
correctness is rarely audited after an annotation is made. We built an agentic review system
in which a large-language-model agent, constrained by a typed schema, a cached literature
corpus and verbatim-quote validation, re-adjudicates every existing GO annotation for a gene
against the primary literature, sequence analysis and curation guidelines, and records a
structured action (\act{ACCEPT}, \act{KEEP\_AS\_NON\_CORE}, \act{MODIFY}, \act{REMOVE},
\act{MARK\_AS\_OVER\_ANNOTATED}, \ldots) with its justification. Across \nReviews{} gene reviews
in \nOrgFolders{} taxa we adjudicated \nAnn{} annotations, citing \nRefs{} references.
\pctChangedAll\% of annotations were changed or rejected (\pctChangedNoPB\% after excluding the
uninformative \emph{protein binding} term). Errors are not uniformly distributed: they
concentrate in specific pipelines (\todo{top pipelines}), specific evidence types (expression
patterns, orthology transfer), and specific ontology regions. We describe ten systematic error
classes in depth---from keyword mappings that conflate regulation with participation to
pseudoenzymes that inherit catalytic activity---several of which have already led to term
obsoletions and pipeline retirements upstream. \todo{one sentence on validation against expert
curators.} Agentic review is a practical route to continuous, auditable quality control of
community knowledge bases.
\end{abstract}
